\documentclass[12pt, a4paper]{article}

% ── Paquetes ────────────────────────────────────────────────────────────────
\usepackage[utf8]{inputenc}
\usepackage[T1]{fontenc}
\usepackage[spanish]{babel}
\usepackage{geometry}
\usepackage{hyperref}
\usepackage{titlesec}
\usepackage{enumitem}
\usepackage{parskip}
\usepackage{fancyhdr}
\usepackage{graphicx}
\usepackage{xcolor}
\usepackage{booktabs}
\usepackage{listings}
\usepackage{lmodern}
\usepackage{microtype}
\usepackage{longtable}
\usepackage{array}
\usepackage{tabularx}

% ── Márgenes ────────────────────────────────────────────────────────────────
\geometry{
  top=2.5cm,
  bottom=2.5cm,
  left=2.8cm,
  right=2.8cm
}

% ── Color corporativo ────────────────────────────────────────────────────────
\definecolor{primary}{RGB}{79, 70, 229}   % indigo
\definecolor{codebg}{RGB}{245, 245, 250}
\definecolor{gray}{RGB}{100, 100, 110}

% ── Hipervínculos ────────────────────────────────────────────────────────────
\hypersetup{
  colorlinks=true,
  linkcolor=primary,
  urlcolor=primary,
  citecolor=primary,
  pdftitle={EventConnect - Documentación del Proyecto},
  pdfauthor={Grupo 11}
}

% ── Cabecera y pie ───────────────────────────────────────────────────────────
\pagestyle{fancy}
\fancyhf{}
\fancyhead[L]{\small\textcolor{gray}{EventConnect — Grupo 11}}
\fancyhead[R]{\small\textcolor{gray}{STW 2025/2026}}
\fancyfoot[C]{\small\textcolor{gray}{\thepage}}
\renewcommand{\headrulewidth}{0.4pt}

% ── Títulos de sección ───────────────────────────────────────────────────────
\titleformat{\section}
  {\large\bfseries\color{primary}}
  {\thesection.}{0.6em}{}[\vspace{-4pt}\rule{\linewidth}{0.6pt}\vspace{2pt}]

\titleformat{\subsection}
  {\normalsize\bfseries}
  {\thesubsection.}{0.5em}{}

% ── Código fuente ────────────────────────────────────────────────────────────
\lstset{
  backgroundcolor=\color{codebg},
  basicstyle=\small\ttfamily,
  breaklines=true,
  frame=single,
  framerule=0pt,
  rulecolor=\color{codebg},
  xleftmargin=8pt,
  xrightmargin=8pt,
  aboveskip=8pt,
  belowskip=8pt
}

% ── Estilo de tablas ─────────────────────────────────────────────────────────
\renewcommand{\arraystretch}{1.3}

% ════════════════════════════════════════════════════════════════════════════
\begin{document}

% ── Portada ──────────────────────────────────────────────────────────────────
\begin{titlepage}
  \centering
  \vspace*{3cm}
  {\Huge\bfseries\color{primary} EventConnect\par}
  \vspace{0.4cm}
  {\large Documentación del Proyecto — Grupo 11\par}
  \vspace{0.8cm}

  \begin{tabular}{ll}
    \textbf{Miembro del grupo:} & Borrel Seral, David (871643) \\[4pt]
    \textbf{Miembro del grupo:} & Báscones Gállego, Pablo (874802) \\[4pt]
    \textbf{Miembro del grupo:} & Hernández Pereda, Mario (873094) \\[4pt]
    \textbf{Miembro del grupo:} & Caudevila Ruiz, Mario (870421) \\
  \end{tabular}

  \vfill
  {\small Sistemas y Tecnologías Web — Curso 2025/2026\par}

  \begin{figure}[h]
      \centering
      \includegraphics[width=0.8\linewidth]{image.png}
      \label{fig:portada}
  \end{figure}
\end{titlepage}

\tableofcontents
\newpage

% ════════════════════════════════════════════════════════════════════════════
\section{URLs de acceso al API y al front-end}

\begin{center}
\begin{tabular}{ll}
  \toprule
  \textbf{Recurso} & \textbf{URL} \\
  \midrule
  Frontend        & \textit{https://frontend-c9bg.onrender.com/} \\
  API / Swagger   & \textit{https://backend-hi6i.onrender.com/api/docs/} \\
  \bottomrule
\end{tabular}
\end{center}

La aplicación ya está desplegada en \textbf{Render} y se puede acceder a ella mediante
los enlaces indicados en la tabla. Se ha desplegado por separado la parte del frontend,
que contiene la interfaz de usuario, y la parte del backend, donde se encuentra la API
y la documentación de Swagger.

Es importante tener en cuenta que, al estar alojada en Render, la primera vez que se
abre alguno de los enlaces puede tardar unos segundos en cargar. Esto ocurre porque el
servicio puede quedarse en reposo cuando no se está usando y necesita unos instantes
para volver a iniciarse. Después de esa primera carga, el funcionamiento suele ser más
rápido y normal.

% ════════════════════════════════════════════════════════════════════════════
\section{Credenciales de acceso}

\begin{center}
\begin{tabular}{lll}
  \toprule
  \textbf{Rol} & \textbf{Email} & \textbf{Contraseña} \\
  \midrule
  Usuario       & \texttt{juan@example.com} & \texttt{password123} \\
  Usuario & \texttt{ana@example.com}   & \texttt{password123} \\
  Usuario & \texttt{carlos@example.com}   & \texttt{password123} \\
  Administrador & \texttt{admin@example.com}   & \texttt{admin123} \\
  \bottomrule
\end{tabular}
\end{center}

% ════════════════════════════════════════════════════════════════════════════
\section{Arquitectura de componentes}

La aplicación sigue una \textbf{arquitectura por capas} con los siguientes componentes principales:

\begin{itemize}[leftmargin=1.4em]
  \item \textbf{Capa de Presentación:} Frontend SPA desarrollado con \textbf{Angular 21}
    (Standalone Components, SSR). Integra \textbf{Leaflet} para mapas interactivos y
    \textbf{Chart.js} para analíticas. Se comunica con el backend mediante HTTP/JSON.

  \item \textbf{Capa de Lógica de Negocio:} Backend API REST con \textbf{Node.js + Express 5}.
    Gestiona autenticación JWT, lógica de eventos, funcionalidades de IA
    y sincronización automática mediante cron job.

  \item \textbf{Capa de Datos:} Base de datos \textbf{MongoDB} gestionada con Mongoose.

  \item \textbf{Servicios Externos:}
    \begin{itemize}
      \item API de Datos Abiertos del Ayuntamiento de Zaragoza (fuente de eventos)
      \item \textbf{Google Gemini 2.5 Flash} para recomendaciones y resúmenes con IA
      \item \textbf{OpenStreetMap} via Leaflet para visualización de eventos en mapa
      \item \textbf{Google OAuth 2.0} para autenticación social
      \item \textbf{Google Maps} para navegación al lugar del evento (enlace directo)
    \end{itemize}
\end{itemize}

\begin{figure}[h]
    \centering
    \includegraphics[width=1\linewidth]{diagrama3.png}
    \caption{Diagrama de componentes}
    \label{fig:diagdeComponentes}
\end{figure}

\subsection{Flujo principal}

\begin{enumerate}[leftmargin=1.6em]
  \item El frontend Angular se comunica con el backend mediante HTTP/JSON.
  \item El backend expone una API REST documentada con Swagger/OpenAPI 3.0.
  \item Un cron job sincroniza eventos cada 6 horas desde la API del Ayuntamiento de Zaragoza.
  \item Los datos se almacenan en MongoDB para mejorar rendimiento y permitir analíticas.
  \item Las funcionalidades de IA se delegan a \textbf{Google Gemini 2.5 Flash}.
\end{enumerate}

% ════════════════════════════════════════════════════════════════════════════
\section{Fuentes de datos abiertos}

\subsection{API del Ayuntamiento de Zaragoza}

El proyecto consume el catálogo de datos abiertos de la Sede Electrónica del Ayuntamiento
de Zaragoza. El recurso principal es el catálogo de Agenda cultural.

\textbf{URL base:}

\begin{lstlisting}
https://www.zaragoza.es/sede/servicio/cultura/evento.json
\end{lstlisting}

\textbf{Parámetros utilizados:}

\begin{center}
\begin{tabular}{ll}
  \toprule
  \textbf{Parámetro} & \textbf{Descripción} \\
  \midrule
  \texttt{start} & Índice de inicio para paginación \\
  \texttt{rows}  & Número de registros por página \\
  \texttt{q}     & Filtros avanzados con operadores OData (\texttt{ge}, \texttt{le}) \\
  \texttt{fl}    & Lista de campos a devolver (reduce el tamaño de la respuesta) \\
  \bottomrule
\end{tabular}
\end{center}

\textbf{Ejemplo de consulta:}

\begin{lstlisting}
https://www.zaragoza.es/sede/servicio/cultura/evento.json
  ?start=0&rows=50
  &q=startDate=ge=2025-01-01T00:00:00Z;endDate=le=2025-12-31T00:00:00Z
\end{lstlisting}

\subsection{Integración en el sistema}

\begin{itemize}[leftmargin=1.4em]
  \item La integración de eventos externos se ha diseñado separando el proceso de
    sincronización con la API del Ayuntamiento de Zaragoza del consumo de los datos
    desde el frontend. De esta forma, la aplicación no consulta directamente la API
    externa cada vez que un usuario accede a los eventos, sino que mantiene una copia
    normalizada en MongoDB.

  \item Un \textbf{cron job} implementado con \texttt{node-cron} se ejecuta cada 6 horas
    y lanza el proceso de sincronización con la API del Ayuntamiento, paginando todos
    los eventos disponibles.

  \item La sincronización también puede ejecutarse manualmente mediante el endpoint
    \texttt{PUT /api/zaragoza/sync}. Además, el endpoint
    \texttt{PUT /api/zaragoza/import} permite importar los eventos obtenidos desde la
    API externa y persistirlos en la base de datos.

  \item Durante la importación, cada evento externo se almacena en MongoDB utilizando
    un campo \texttt{externalId} único. Este identificador permite detectar si el evento
    ya existe, evitando duplicados y permitiendo actualizar la información cuando el
    evento vuelve a recibirse desde la API externa.

  \item Los eventos caducados no se eliminan inmediatamente, sino que se marcan con
    \texttt{status: 'expired'}. Esto permite conservar histórico para estadísticas y
    mantener la trazabilidad de eventos ya finalizados.
    Los eventos que permanecen más de 12 meses en estado \texttt{expired} se
    eliminan automáticamente, evitando un crecimiento indefinido de la base de datos.

  \item Si un evento marcado como \texttt{expired} vuelve a aparecer en una
    sincronización posterior, el sistema lo actualiza y puede reactivarlo, cambiando
    su estado de nuevo a activo. Este comportamiento resulta útil para eventos
    recurrentes o para eventos cuya información ha sido actualizada.

  \item El frontend no consume directamente la API del Ayuntamiento, sino los datos ya
    almacenados en MongoDB mediante los endpoints propios de la aplicación. Por
    ejemplo, \texttt{GET /api/events} permite obtener el listado general de eventos,
    \texttt{GET /api/events/:id} permite consultar un evento concreto y
    \texttt{GET /api/events/sections} devuelve secciones preparadas para la página
    principal, como eventos destacados, eventos de hoy, eventos de la semana o eventos
    recientes.

  \item El endpoint \texttt{GET /api/zaragoza} permite consultar eventos directamente
    desde la fuente externa aplicando parámetros como paginación, búsqueda, filtro por
    eventos del día o consulta de un evento concreto. Este endpoint resulta útil para
    pruebas, depuración o consulta puntual de la API externa, mientras que el flujo
    habitual del frontend utiliza los endpoints internos bajo \texttt{/api/events}.

  \item El campo \texttt{geometry} de cada evento se integra con \textbf{OpenStreetMap}
    mediante \textbf{Leaflet}, permitiendo mostrar la ubicación del evento en un mapa
    interactivo dentro de la interfaz.
\end{itemize}

% ════════════════════════════════════════════════════════════════════════════
\section{Módulos del back-end}

A continuación se listan los módulos npm principales utilizados en el backend
junto con una breve descripción de su función.

\subsection{Dependencias del proyecto}

\begin{center}
\begin{longtable}{p{4.2cm} p{1.8cm} p{7.5cm}}
  \toprule
  \textbf{Módulo} & \textbf{Versión} & \textbf{Descripción} \\
  \midrule
  \endhead
  \texttt{express}              & \^{}5.2.1  & Framework web para la API REST \\
  \texttt{mongoose}             & \^{}9.3.0  & ODM para MongoDB; define modelos y esquemas \\
  \texttt{@google/generative-ai}& \^{}0.24.1 & SDK oficial de Google para Gemini; usado en recomendaciones e IA \\
  \texttt{node-cron}            & \^{}4.2.1  & Programación de tareas periódicas (sincronización de eventos cada 6h) \\
  \texttt{axios}                & \^{}1.13.6 & Cliente HTTP para consumir la API del Ayuntamiento \\
  \texttt{jsonwebtoken}         & \^{}9.0.3  & Generación y verificación de tokens JWT para autenticación \\
  \texttt{bcrypt / bcryptjs}    & \^{}6.0.0  & Hash seguro de contraseñas de usuarios \\
  \texttt{google-auth-library}  & \^{}10.6.1 & Verificación de tokens de Google OAuth 2.0 \\
  \texttt{express-validator}    & \^{}7.3.1  & Validación y saneamiento de entradas en las rutas \\
  \texttt{swagger-jsdoc}        & \^{}6.2.8  & Generación automática de especificación OpenAPI desde JSDoc \\
  \texttt{swagger-ui-express}   & \^{}5.0.1  & Interfaz visual interactiva Swagger UI \\
  \texttt{helmet}               & \^{}8.1.0  & Cabeceras de seguridad HTTP \\
  \texttt{cors}                 & \^{}2.8.6  & Gestión de política de origen cruzado (CORS) \\
  \texttt{dotenv}               & \^{}17.3.1 & Carga de variables de entorno desde \texttt{.env} \\
  \texttt{morgan}               & \^{}1.10.1 & Logger de peticiones HTTP \\
  \texttt{nodemailer}           & \^{}8.0.2  & Envío de correos electrónicos (notificaciones) \\
  \texttt{proj4}                & \^{}2.20.4 & Conversión de coordenadas geográficas \\
  \texttt{cookie-parser}        & \^{}1.4.7  & Parseo de cookies en las peticiones \\
  \texttt{ajv / ajv-formats}    & \^{}8.18.0 & Validación de esquemas JSON \\
  \texttt{jest}                  & \^{}30.3.0  & Framework de tests unitarios e integración \\
  \texttt{supertest}             & \^{}7.2.2   & Tests de integración para la API HTTP \\
  \texttt{mongodb-memory-server} & \^{}11.0.1  & Servidor MongoDB en memoria para tests (sin BD real) \\
  \texttt{nodemon}               & \^{}3.1.14  & Recarga automática del servidor en desarrollo \\
  \bottomrule
\end{longtable}
\end{center}

% ════════════════════════════════════════════════════════════════════════════
\section{Prototipado de la solución}

El prototipo completo de la aplicación está disponible en Figma:

\begin{center}
  \textbf{\textcolor{primary}{\href{https://www.figma.com/design/TfRqhXv41VX0Toq1gj84KK/EventConnect?node-id=0-1&t=e7v9mcGCSKsCXTvy-1}{Prototipo de la versión web en Figma}}}
\end{center}

\begin{center}
  \textbf{\textcolor{primary}{\href{https://www.figma.com/design/v0OXOjUJ2Dfx8RJBapv4A1/EventConnect-movil?node-id=1-2&t=tOttEwtRHJ7NRy7p-0}{Prototipo de la versión móvil en Figma}}}
\end{center}

El prototipo cubre todas las pantallas de la aplicación:
\begin{itemize}[leftmargin=1.4em]
  \item Exploración y detalle de eventos
  \item Visualización en mapa interactivo
  \item Perfil de usuario e historial
  \item Chat de evento y mensajería privada
  \item Sistema de quedadas y amigos
  \item Panel de analíticas
  \item Panel de administración
\end{itemize}

% ════════════════════════════════════════════════════════════════════════════
\section{Front-end: tecnología y módulos}

\textbf{Framework:} Angular 21 con Standalone Components y Server-Side Rendering
(\texttt{@angular/ssr}).

\begin{center}
\begin{longtable}{p{4.2cm} p{1.8cm} p{7.5cm}}
  \toprule
  \textbf{Módulo} & \textbf{Versión} & \textbf{Descripción} \\
  \midrule
  \endhead
  \texttt{@angular/core}            & \^{}21.2.0 & Núcleo del framework Angular \\
  \texttt{@angular/router}          & \^{}21.2.0 & Enrutamiento SPA con lazy loading \\
  \texttt{@angular/forms}           & \^{}21.2.0 & Formularios reactivos y validación \\
  \texttt{@angular/common}          & \^{}21.2.0 & Directivas, pipes y servicios comunes \\
  \texttt{@a/platform-browser}& \^{}21.2.0 & Renderizado en navegador \\
  \texttt{@a/platform-server} & \^{}21.2.0 & Renderizado en servidor (SSR) \\
  \texttt{@angular/ssr}             & \^{}21.2.1 & Server-Side Rendering para Angular \\
  \texttt{leaflet}                  & \^{}1.9.4  & Mapas interactivos con OpenStreetMap para visualización de eventos \\
  \texttt{chart.js}                 & \^{}4.5.1  & Gráficas de analíticas (barras, líneas, popularidad) \\
  \texttt{rxjs}                     & \textasciitilde7.8.0 & Programación reactiva y gestión de peticiones asíncronas \\
  \texttt{proj4}                    & \^{}2.20.4 & Conversión de coordenadas geográficas para el mapa \\
  \texttt{vitest}                   & \^{}4.0.8  & Framework de tests unitarios para el frontend \\
  \texttt{jsdom}                    & \^{}28.0.0 & Entorno DOM para tests \\
  \texttt{prettier}                 & \^{}3.8.1  & Formateo automático de código \\
  \bottomrule
\end{longtable}
\end{center}

% ════════════════════════════════════════════════════════════════════════════
\section{Modelo de IA}

\subsection{Modelo utilizado: Google Gemini 2.5 Flash Lite}

Se utiliza \textbf{Gemini 2.5 Flash Lite} de Google AI como motor de inteligencia
artificial. La justificación de esta elección es su equilibrio óptimo entre
\textbf{coste, velocidad de respuesta y calidad} para tareas de texto:

\begin{itemize}[leftmargin=1.4em]
  \item Al tratarse de una aplicación orientada al consumidor con llamadas frecuentes a
    la IA, los modelos más potentes (Gemini Pro, GPT-4) resultarían en costes demasiados elevados, y no sería eficiente puesto que se dispondría de una capacidad de computación mucho mayor de la necesaria.
  \item Flash Lite ofrece respuestas rápidas y coherentes para las funcionalidades
    implementadas (clasificación de intención y generación de texto estructurado).
  \item Permite monitorizar y controlar el consumo desde Google AI Studio.
\end{itemize}

\textbf{Control de costes implementado:}
\begin{itemize}[leftmargin=1.4em]
  \item Rate limiting en el backend para limitar llamadas a la API de Gemini.
  \item Uso del modelo \texttt{gemini-2.5-flash-lite} (tier más económico).
  \item Monitorización del consumo desde Google AI Studio.
\end{itemize}

\subsection{Recomendador SpAlk (\texttt{POST /api/recommend})}

\textbf{Descripción:} Sistema conversacional que interpreta la intención del usuario
(con quién va y qué tipo de plan busca) para recomendar eventos reales de la base de datos.
Utiliza un mapa semántico de categorías para traducir la intención a filtros reales sobre
MongoDB.

\textbf{Prompt utilizado:}

\begin{lstlisting}
Eres SpAlk, un asistente experto en recomendar planes en Zaragoza.

Tu funcion NO es inventar eventos.
Tu funcion es interpretar la intencion del usuario para filtrar
eventos reales en una base de datos.

El usuario quiere:
- Ir: {companionText}
- Busca: {vibeText}

Tu tarea:
1. Analiza la intencion
2. Devuelve SOLO una categoria de evento de esta lista:
   deporte / musica / cultura / gastronomia / social / naturaleza

Reglas:
- Responde SOLO con una palabra
- Sin frases, sin explicaciones, sin texto adicional
\end{lstlisting}


\subsection{Resumen inteligente (\texttt{POST /api/ai/summary})}

\textbf{Descripción:} Genera un resumen de los eventos disponibles
(filtrables por categoría y rango de fechas), identificando tendencias y destacando
los más relevantes.

\textbf{Prompt utilizado:}

\begin{lstlisting}
Eres un sistema que analiza eventos en Zaragoza.

Tu tarea tiene DOS PARTES:

1. Generar un RESUMEN general (5 a 10 lineas):
   - Describe tendencias (categorias, tipos de eventos)
   - NO menciones nombres concretos
   - NO recomiendes

2. Seleccionar 3 a 5 eventos destacados

FORMATO (JSON PURO, sin markdown, sin bloques de codigo):
{
  "summary": "texto con saltos de linea",
  "highlights": [
    { "text": "breve descripcion", "eventId": "id" }
  ]
}

EVENTOS:
{JSON con los primeros 20 eventos: id, titulo, categoria, fecha}
\end{lstlisting}

\subsection{Conexión entre la IA y la base de datos}

La IA no se conecta directamente a MongoDB ni tiene acceso propio a la base de datos.
La conexión se realiza a través del backend de la aplicación. Es decir, cuando el
frontend necesita una recomendación o un resumen, envía una petición al backend, y es
el backend quien consulta MongoDB y le envía la información necesaria a Gemini.

En el caso del recomendador, Gemini no recibe todos los eventos de la base de datos.
Solo recibe la intención del usuario, por ejemplo con quién quiere ir y qué tipo de
plan busca. A partir de eso, la IA devuelve una categoría. Después, el backend usa esa categoría para buscar en MongoDB
eventos reales que coincidan con ella.

En el caso del resumen inteligente, el proceso es algo distinto. Primero el backend
consulta MongoDB para obtener una lista de eventos activos, pudiendo aplicar filtros
como categoría o fecha. Después, esos eventos se envían a Gemini para
que genere un texto explicando qué tipo de eventos hay y cuáles pueden destacar. De
esta forma, la IA puede ``saber'' qué eventos existen, pero solo porque el backend se
los proporciona previamente.

Por tanto, MongoDB sigue siendo la fuente real de información de la aplicación. Gemini
se utiliza como una capa de apoyo para interpretar la información así como la intención del usuario

% ════════════════════════════════════════════════════════════════════════════
\section{Validación y pruebas}

\subsection{Tests unitarios e integración (backend)}

Se utilizan \textbf{Jest} + \textbf{Supertest} + \textbf{mongodb-memory-server}:

\begin{itemize}[leftmargin=1.4em]
  \item Tests sobre los controladores de la API: respuestas HTTP, códigos de estado,
    validaciones de entrada y manejo de errores.
  \item Tests sobre los modelos Mongoose: esquemas, validaciones y métodos.
  \item \texttt{mongodb-memory-server} para levantar una BD en memoria, sin dependencia
    de instancia real en el entorno de CI.
\end{itemize}

La cobertura de tests se ejecuta con:

\begin{lstlisting}[language=bash]
cd backend && npm test -- --coverage
\end{lstlisting}

\subsection{Tests End-to-End (frontend)}

Se han implementado \textbf{2 pruebas E2E con Playwright}:

\textbf{Test 1 — Login de usuario:}
Verifica el flujo completo de autenticación: acceso a la página de login, introducción
de credenciales de usuario válidas y redirección correcta al dashboard tras la
autenticación exitosa.

\textbf{Test 2 — Apuntarse a un evento:}
Verifica que un usuario autenticado puede navegar al detalle de un evento, pulsar el
botón de inscripción y que la aplicación confirma correctamente su participación en el evento.

% ════════════════════════════════════════════════════════════════════════════
\section{Mejoras implementadas}

\subsection{CI/CD con GitHub Actions}

Se ha configurado un pipeline de integración continua que se ejecuta automáticamente
en cada \textit{push} y \textit{pull request} a la rama \texttt{main}:

\begin{itemize}[leftmargin=1.4em]
  \item Instalación de dependencias del backend.
  \item Ejecución del conjunto completo de tests con Jest.
  \item Verificación de que el build del frontend Angular compila correctamente.
\end{itemize}

Esto garantiza que ningún cambio rompe los tests ni la compilación antes de fusionarse
con la rama principal, facilitando la detección de posibles problemas.

\subsection{Cobertura de tests}

Se ha integrado la generación de informes de cobertura de código en el pipeline de CI.
Esto proporciona visibilidad continua sobre qué partes del backend están cubiertas por
tests.

% ════════════════════════════════════════════════════════════════════════════
\section{Valoración global del proyecto}

EventConnect ha sido un proyecto ambicioso que nos ha permitido aplicar de forma real
el stack MEAN completo, integrando además tecnologías externas como IA generativa, datos
abiertos del Ayuntamiento de Zaragoza y mapas interactivos.

El punto más satisfactorio ha sido ver cómo los distintos módulos (sincronización
automática de eventos, recomendador con IA, chat en tiempo real, sistema social) se
integran en una aplicación funcional
que resuelve una necesidad real: conectar a las personas con los eventos culturales y
de ocio de su ciudad.



% ════════════════════════════════════════════════════════════════════════════
\section{Mejoras propuestas}

Si volviéramos a empezar o tuviéramos más tiempo, identificamos las siguientes
mejoras y limitaciones conocidas:

\begin{itemize}[leftmargin=1.4em]
  \item \textbf{WebSockets reales para el chat:} la mensajería en tiempo real actual
    podría beneficiarse de una implementación más robusta con Socket.io y salas
    dedicadas por evento.
  \item \textbf{Caché de eventos:} implementar Redis para cachear los resultados de la
    API del Ayuntamiento y reducir latencia en las consultas más frecuentes.
  \item \textbf{Notificaciones push:} integrar Web Push para recordatorios automáticos
    de eventos a los que el usuario se ha inscrito.
  \item \textbf{Búsqueda semántica:} usar \textit{embeddings} con la IA para búsqueda
    por similitud semántica en lugar de solo filtros por categoría exacta.
  \item \textbf{Limitación conocida:} el sistema de recomendaciones depende de que los
    eventos de la API del Ayuntamiento estén correctamente categorizados. Si hay cambios en las categorías podría haber problemas en el caso de que aparezcan nuevas o desaparezcan las actuales.
\end{itemize}

\end{document}
